Suponga que un usuario genera al azar dos números primos $P$ y $Q$ distintos de 600 dígitos cada uno, y con esto define $N = P \cdot Q$. Suponga además que este usuario elige como clave pública $e = 3$, asegurándose de que $e$ sea primo relativo con $\phi(N)$. Suponga que un atacante sabe que se enviará el texto cifrado de algún mensaje $m$ tal que $m < \sqrt[3]{N}$. Demuestre que el atacante puede descifrar dicho texto cifrado en tiempo polinomial.

\paragraph{Solución.} Dado que $m < \sqrt[3]{N}$, se tiene que $m^3 < N$. Tenemos entonces que el mensaje cifrado que será enviado en este caso es $c = m^3 \,\text{mod}\, N = m^3$, puesto que la clave pública es $e = 3$. De esta forma, un atacante puede calcular $m$ a partir de $c$ simplemente calculando la raíz cúbica de $c$, lo cuál se puede hacer de manera eficiente utilizando búsqueda binaria.
